\documentclass{article}
\input{../../../../lib.sty}

\title{\texttt{fn match\_per\_group\_predicate}}
\author{\MichaelShoemateAuthor}
\begin{document}
\maketitle


\contrib

Proves soundness of \rustdoc{transformations/make_stable_lazyframe/truncate/matching/fn}{match\_per\_group\_predicate}
in \asOfCommit{mod.rs}{f5bb719}.

\section{Hoare Triple}
\subsection*{Precondition}
\subsubsection*{Compiler-verified}
Types matching pseudocode.

\subsubsection*{Caller-verified}
The threshold is normalized to the maximum number of retained ranks or indices.

\subsection*{Pseudocode}
\label{sec:python-pseudocode}
\lstinputlisting[language=Python,firstline=2,escapechar=|]{./pseudocode/match_per_group_predicate.py}

\subsection*{Postcondition}
\begin{theorem}[Postcondition]
    If \texttt{enumeration} is a trusted uniform permutation of row indices
    or ordinal ranks of sort keys with trusted uniform tie-breaking,
    and \texttt{partition\_by} includes \texttt{identifier},
    then returns a candidate per-group contribution bound together with its sort keys.
    Validating these keys for permutation equivariance and supported scalar dtypes
    makes the bound valid for the remaining grouping expressions in \texttt{partition\_by}.
\end{theorem}

\begin{proof}
    The matcher accepts only a full trusted uniform permutation of
    \texttt{int\_range(0, len)}, or ordinal ranks of a struct consisting
    of sort keys and a final trusted uniform permutation.
    Line \ref{line:check_enumeration} rejects other expressions.

    In the sorted case, each row has a distinct random tie-breaking priority.
    The ordinal-rank matcher on line \ref{line:check_enumeration} returns all sort keys
    as deferred validation obligations.
    The caller checks them for permutation equivariance, infallibility, non-resizing
    behavior, and supported scalar dtypes against the preceding plan's output
    schema before using the bound for stability accounting. Arbitrary nested
    dtypes and physical categorical encodings are excluded from ranking.

    The enclosing window evaluates the ranking within each partition,
    which contains the identifier. Consequently, changing one identifier's
    rows cannot affect keys in another identifier's partition. The uniform
    tie-breaker makes selection invariant in distribution under physical
    input-row permutations.

    By line \ref{line:partition}, the window partitions rows by identifier
    and the additional grouping expressions. Full permutations assign
    distinct values $0,\ldots,n-1$. Ordinal ranks assign $1,\ldots,n$.
    The caller normalizes the threshold to the number of retained indices or ranks.
    Line \ref{line:bound} records that bound.
\end{proof}
\end{document}
